% Generated draft: requires mathematical and source review.
\documentclass[11pt]{article}
\usepackage[margin=25mm]{geometry}
\usepackage{fontspec}
\setmainfont{Latin Modern Roman}
\usepackage{amsmath,amssymb,mathtools,unicode-math}
\setmathfont{Latin Modern Math}
\usepackage{enumitem,needspace,xurl,hyperref}
\setlength{\parindent}{0pt}
\setlength{\parskip}{5pt}
\setlength{\emergencystretch}{3em}
\newcommand{\N}{\mathbb{N}}
\newcommand{\Z}{\mathbb{Z}}
\newcommand{\Q}{\mathbb{Q}}
\newcommand{\R}{\mathbb{R}}
\newcommand{\C}{\mathbb{C}}
\newcommand{\F}{\mathbb{F}}
\newcommand{\A}{\mathbb{A}}
\newcommand{\PP}{\mathbb{P}}
\newcommand{\E}{\mathbb{E}}
\newcommand{\eps}{\varepsilon}
\newcommand{\dd}{\,\mathrm{d}}
\begin{document}
\section*{UnsolvedMath ID: 30000493}
\textbf{Source problem code:} catalogue entry
\subsection{Definitions and mathematical statement}
The restored record contains three questions: dense settling for quadratic arboreal images, an explicit pro-$2$ kernel, and a natural tree zeta/$L$-function.

In the usual basepoint-zero construction, vertices at level $n$ are roots of $f^{\circ n}(x)=0$, with edges given by $f$ and root $0$. When all levels are distinct and have the expected branching, this is a regular rooted binary tree $T_f$. Galois acts on it, giving
\[
 \rho_f:\operatorname{Gal}(\overline\Q/\Q)\to\operatorname{Aut}(T_f),
 \qquad G_f=\operatorname{im}\rho_f.
\]
The basepoint and noncollision hypotheses are not stated in the excerpt. For arbitrary quadratic $f\in\Z[x]$, repeated preimages or a basepoint in the postcritical orbit can invalidate this regular-tree setting; those cases must not be silently included in a regular-binary-tree theorem or silently discarded.

For a regular binary tree, a cycle of $g$ on level $n$ is stable if, at every deeper level, all descendants of that cycle form one cycle. Call $g$ settled if the fraction of level-$n$ vertices in stable cycles tends to one as $n\to\infty$. A closed subgroup $G\leq\operatorname{Aut}(T_f)$ is densely settled if its settled elements are dense in $G$: for every $g\in G$ and finite level there is a settled element of $G$ with the same action through that level. This is neither density of $G$ in the full tree automorphism group nor a Haar-almost-everywhere assertion. The first question asks whether the quadratic-integral arboreal groups in the intended admissible setting have this property.

For the explicit polynomial
\[
 f(x)=(x+1)^2-2,
\]
let $K_\infty$ be the union of splitting fields of its iterates in the intended construction, and let $\Q_{\{2,\infty\}}(2)$ be the maximal pro-$2$ extension of $\Q$ unramified outside $\{2,\infty\}$. The source's quotient presupposes the corresponding inclusion of fields, and asks to determine
\[
 \ker\!\left(\operatorname{Gal}(\Q_{\{2,\infty\}}(2)/\Q)
      \twoheadrightarrow G_f\right)
 =\operatorname{Gal}(\Q_{\{2,\infty\}}(2)/K_\infty).
\]
The original asks its size, while the corrected statement says determine the kernel; whether this means pro-$2$ structure, generator rank or another size invariant is not specified.

Finally one is to construct a natural zeta function or $L$-function for the rooted tree of iterates. The arithmetic data to encode, normalization, finite-level compatibility, Euler factors and convergence are not prescribed; no arbitrary Euler product is substituted.

The supplied literature assessment checked 21 August 2026 reports densely settled groups in important families, including special postcritically finite cases, but not a general solution of the quadratic-integral and kernel questions. The separate recent special-case reference and all three tasks are retained with that dated qualification.
\subsection{Short English statement}
Are the relevant quadratic arboreal Galois groups densely settled? Determine the stated pro-$2$ kernel for $(x+1)^2-2$, and construct a natural zeta or $L$-function for the iterate tree. Basepoint/regular-tree hypotheses, the kernel-size target and the zeta data remain to be specified; supplied special families do not settle every part.
\subsection{Sources}
\begin{itemize}
\item Local UnsolvedMath registry, permanent ID 30000493, exact source code catalogue entry; \emph{Dense Settling and Arboreal Galois Representations}. Proposed year 2006, status partially\_solved. The corrected\_verified statement restores the explicit polynomial, ramification set and three questions; its claimed inspection is source metadata.
\item Supplied report citation: \emph{Pro-$p$ Extensions of Global Fields and pro-$p$ Groups (2007)}, Oberwolfach Reports, DOI 10.4171/owr/2006/25, pp. 1463--1536, \url{https://doi.org/10.4171/owr/2006/25}. Title year differs from DOI/proposed year.
\item Supplied literature triage checked 2026-08-21: ``original OWR report'', \url{https://doi.org/10.4171/owr/2006/25}; ``Boston--Jones primary paper'', \url{https://people.carleton.edu/~rfjones/Preprints/SerreImageNew2.pdf}; ``recent special-case result'', \url{https://arxiv.org/abs/2604.04524}. The dated OWR-derived curation requests independent verification of time-sensitive status and cited claims; no linked PDF was accessed here.
\end{itemize}
\end{document}
